
\subsection{Theorie} \label{sec:theorie}

A partir de Ref. \cite{protocoleDMA}, nous pouvons decrire depuis les lois de Hook et de Newton la modelisation de notre systeme. L'instrument imporse une contrainte périodique en fonction du temps $\sigma$(t) de pusation $\omega$ et d'une contrainte de reference $\sigma_0$ tel que 
\begin{equation}
    \sigma(t) = \sigma_0 \sin(\omega t).
	\label{eq:1}
\end{equation}

Notons, une deformation $\epsilon$ sous dephasage d’un angle $\delta$ :
\begin{equation}
    \varepsilon(t) = \varepsilon_0 \sin(\omega t + \delta).
	\label{eq:2}
\end{equation}

Selon les hypotheses du protocole (Ref. \cite{protocoleDMA} p. 2-3), les grandeurs viscoelastiques mesurees $E'$ (module de stockage) et $E''$ (module de perte) sont apres derivations des Eqs. \ref{eq:1} et \ref{eq:2}
\begin{equation}
	E' &= \frac{\sigma_0}{\varepsilon_0}\cos\delta \hspace{0.4cm} et \hspace{0.4cm}
    E'' &= \frac{\sigma_0}{\varepsilon_0}\sin\delta.
	\label{eq:3}
\end{equation}

Elles permettent d'ecrire le pic du facteur d’amortissement $\tan\delta$ correspondant à la transition vitreuse mécanique a partir de relations trigonometriques avec Eqs. \ref{eq:3} tel que :
\begin{equation}
	\tan\delta &= \frac{E'}{E'\'}.
	\label{eq:4}
\end{equation}

Comme cette transition depend de la frequence d’excitation, les temperatures des pics mesures
sont exploitees à l’aide de la loi d’Arrhenius (voir Ref. \cite{protocoleDMA}) :
\begin{equation}
    f = f_0 \exp\!\left(-\frac{E_a}{RT}\right),
	\label{eq:arrhenius}
\end{equation}

que l’on linéarise sous la forme

\begin{equation}
    \ln f = \ln f_0 - \frac{E_a}{R}\,\frac{1}{T}.
	\label{eq:arrheniuslin}
\end{equation}

Ainsi, la représentation de $\ln f$ en fonction de $1/T_g(f)$ permet d’extraire l’énergie
d’activation $E_a$ associée à la transition vitreuse du PMMA.

